\section{Rationale of the study}

This study examines whether the value relevance of accounting information
differed before and after the recent commercial banking consolidation wave in
Nepal. Its contribution is threefold.

First, the Nepalese price-determinant studies reviewed above estimate average
price associations, while merger-performance studies examine bank outcomes
around consolidation, including through recent comparative designs. The
studies reviewed do not test whether EPS--price and BVPS--price slopes changed
across the early-2023 consolidation period. The interaction model addresses
this narrower question for the selected bank series. This is a specific
contribution to the reviewed literature, not a claim that no related research
exists.

Second, the study applies a standardized price-measurement rule on or after the
Chaitra-end target, which operationalizes the statutory nine-month reporting
deadline in Section 59(2) of the Banks and Financial Institutions Act
\citep{bafia2017}. This makes the observation window reproducible across
bank-years. The deadline is a legal requirement, not evidence of timely
publication by each bank; the study will record release dates where verifiable
and will not describe these prices as confirmed post-disclosure observations.

Third, the study adopts a parsimonious model grounded in the residual-income
valuation framework of \citet{ohlson1995}, using EPS and BVPS as the primary
explanatory variables and excluding price-derived ratios such as P/E. This
avoids the definitional overlap that arises when a ratio constructed from the
dependent variable is used as an explanatory variable.

The design does not estimate causal merger effects or support generalization
to all commercial banks. Because the periods are defined by calendar time, the
interaction coefficients capture differences in estimated associations across
periods that may reflect consolidation, macroeconomic conditions, or other
concurrent changes. The small purposive sample, changes in reporting
boundaries, and exploratory inference limit the strength and reach of any
conclusions.
